\section{Discussion and Limitations}
\label{sec:limitations}

The central difficulty is connecting a learned proposal to its actual consequence. An endpoint-conditioned latent path can be plausible under the trajectory model yet difficult to execute. Conversely, a decoder can be accurate on recorded paths and unreliable on generated paths. Evaluating inverse dynamics on real and generated paths separately helps distinguish these cases, but B's learned score is itself imperfect. Sharing its parameters with the generators can introduce correlated errors, so parameter reuse should be evaluated rather than assumed to improve reliability.

The completed A/B study also limits the interpretation of positive transfer. Action dimension, task dynamics, representation quality, and goal difficulty vary together across the four tasks. Cube's larger action block and positive action-supervision effect are therefore a correlation, not evidence that action dimension causes the improvement. Increasing model capacity or switching to recursive prediction has not provided a uniform remedy in the tested follow-ups.

Our current tasks use short goal offsets and finite evaluation cohorts. High actor success on some tasks restricts the room for useful additional planning. Stronger claims about planning should be tested with suitable goal distances, recovery conditions, or longer-horizon tasks under a separately frozen protocol. Current data-exposure uncertainty also prevents a claim of unseen-data generalization for the exploratory four-mode run.

Online adaptation is a separate ongoing study. It will compare frozen dynamics, matched offline continuation, and updates from executed transitions while freezing the visual representation. No positive online-learning or continual-learning result is claimed here. Combining online updates with the four-mode model would introduce additional questions about shared-parameter drift and requires its own evaluation.

\section{Conclusion}

Fast-LeWAM provides a shared latent architecture for action generation, causal future prediction, path generation, and action decoding. It supports a direct comparison of action and latent-path proposals using the same learned consequence scorer. Completed action--dynamics experiments show that joint supervision and actor-initialized planning have task-dependent effects, motivating explicit evaluation of proposal coverage and selection quality. Final four-mode experiments will determine whether the additional path capability yields a useful performance--cost trade-off. The current evidence supports the problem formulation and implemented framework, while the effectiveness of the proposed extension remains open.
